\uuid{loFG}
\titre{Vrai/faux}
\auteur{bodin ; hanani ; mzari}
\organisation{exo7}
\licence{CC BY-NC-SA 4.0}
\datecreate{2020-01-20}
\niveau{L1}
\module{Algèbre}
\chapitre{Logique, ensemble, raisonnement}
\sousChapitre{Logique}
\difficulte{2}
\theme{Logique, vrai, faux}
\exo7qcmid{lille-1.logique.007}

\contenu{
\begin{question}
\qclassification{100.01}
\qtags{motcle=vrai/faux}
Soit $P$ une assertion fausse, $Q$ une assertion vraie et $R$ une assertion fausse. Quelles sont les assertions vraies ?
\begin{answers}
    \bad{$Q$ et ($P$ ou $R$)}

    \bad{$P$ ou ($Q$ et $R$)}

    \good{non($P$ et $Q$ et $R$)}

    \good{($P$ ou $Q$) et ($Q$ ou $R$)}
\end{answers}
\begin{explanations}
Il suffit de remplacer $P$ par "faux", $Q$ par "vrai" et $R$ par "faux". Par exemple "$Q$ et ($P$ ou $R$)" devient "vrai et (faux ou faux)", qui est la même chose que "vrai et faux", qui est donc "faux".
\end{explanations}
\end{question}
}
